\section*{ID7327: Quick Dimensional Check \& Numerics: 10\textasciicircum{}121}
\paragraph{Metadata.}
PiID: 3435139844253 | RTEID: RTE:P28452.3682:T2.8985 | UUID: 544fc0c1-680c-58df-bed7-e340d97ae3da

\paragraph{Canonical Statement.}
So to get the observed value you need a **huge dimensionless suppression or enhancement factor**, i.e. your \textbackslash{}(\textbackslash{}Phi\_0/\textbackslash{}Phi\_\{\textbackslash{}rm total\}\textbackslash{}) must be of order \textbackslash{}(10\textasciicircum{}\{121\}\textbackslash{}) if \textbackslash{}(R\textbackslash{}) is the Hubble radius, or of order \textbackslash{}(10\textasciicircum{}\{-122\}\textbackslash{}) if \textbackslash{}(R\textbackslash{}) is Planck length. This is not a show-stopper, but it must be acknowledged and justified conceptually.

\paragraph{Canonical Equation.}
\[
10^{121}
\]

\paragraph{Dictionary.}
Quick Dimensional Check \& Numerics: 10\textasciicircum{}121 — 10\textasciicircum{}121

\paragraph{Encyclopedia.}
Quick Dimensional Check \& Numerics: 10\textasciicircum{}121 is a symbolic expression, written 10\textasciicircum{}121. It introduces a symbol or relation used elsewhere in the framework. Within the Trawin Zero Point registry it serves as one element of the framework's mathematical narrative.
